\documentclass[11pt]{amsart}
\usepackage{amsmath,amssymb,amsthm,mathtools}
\usepackage[margin=1in]{geometry}
\usepackage{booktabs}
\usepackage{longtable}
\usepackage{hyperref}
\usepackage{graphicx}
\usepackage{xcolor}
\usepackage{enumitem}
\usepackage{tikz-cd}
\hypersetup{colorlinks=true,linkcolor=blue,urlcolor=blue,citecolor=blue}

\newtheorem{theorem}{Theorem}[section]
\newtheorem{conjecture}[theorem]{Conjecture}
\newtheorem{proposition}[theorem]{Proposition}
\newtheorem{lemma}[theorem]{Lemma}
\newtheorem{corollary}[theorem]{Corollary}
\theoremstyle{definition}
\newtheorem{definition}[theorem]{Definition}
\theoremstyle{remark}
\newtheorem{remark}[theorem]{Remark}

\newcommand{\Om}{\widetilde{\Omega}}
\newcommand{\ve}{\varepsilon}
\newcommand{\loc}{\mathrm{loc}}
\newcommand{\PROVED}{\textbf{(PROVED)}}
\newcommand{\COMP}{\textbf{(COMPUTED, PROOF PENDING)}}
\newcommand{\NOGO}{\textbf{(NO-GO)}}
\newcommand{\OPEN}{\textbf{(OPEN)}}
\newcommand{\WALL}{\textbf{(WALL = RH)}}

\title[A Spectral Reduction of RH via Zeta Spectral Triples]{A Spectral Reduction of the
Riemann Hypothesis via Zeta Spectral Triples\\[2pt]
\large with a full record of the sixty-phase program that led to it}
\author{David Alejandro Trejo Pizzo}
\thanks{Independent Researcher, in collaboration with the Connes--Consani--Moscovici program.
\texttt{dtrejopizzo@gmail.com}}
\date{June 2026}

\begin{document}
\maketitle

\begin{abstract}
\noindent We do not claim a proof of the Riemann Hypothesis (RH). What we do is \emph{localize}
it: we reduce RH to a single explicit, phase-sensitive positivity statement, we prove every
surrounding link in the chain that leads to that statement, we record every dead end of a
sixty-phase program together with its precise obstruction, and we build --- on the
Connes--Consani--Moscovici (CCM) framework of zeta spectral triples --- the prolate-operator
scaffolding that surrounds the residual wall. The architecture is a ladder of twelve steps:
the first ten are proved or constructed-and-validated, they pin the truth of RH to the
positivity of the ground-state eigenvalue $\ve_0(\lambda)$ of a localized Weil quadratic form
$A_\lambda$, and the last two identify that positivity as the one genuinely RH-hard residue ---
\emph{the wall} --- reached identically from primes, zeros, the prolate operator, and de
Branges theory. A novel contribution of the present version is the isolation of a single new
lemma (the CCM shorted Dirichlet quotient, Lemma~\ref{lem:23F}) on which the uniformity of the
structure rests; we show it is a genuinely new theorem about the Weil form, not a consequence
of the existing prolate literature, and we report how far its proof has been carried. All
statements carry a status mark: \PROVED{} (rigorous proof or citation that holds); \COMP{}
(high precision, proof pending); \NOGO{} (tried and refuted, with reason); \OPEN{} (open but
not RH-hard); \WALL{} (logically equivalent to RH).
\end{abstract}

\tableofcontents
\newpage

\section{Introduction: the program from Phase 1 to Phase 60}

The aim has never been to \emph{prove} RH by frontal assault, but to \emph{localize} it: to
reduce RH to a single, explicit, finite, phase-sensitive statement, with every surrounding link
proved, every dead end recorded with its obstruction, and the genuinely hard residue named.
Sixty phases produced (i) several rigorous unconditional results, (ii) a complete catalogue of
what does not work and why, (iii) a precise identification --- from seven independent
directions --- of a single master wall, and (iv) the present reduction onto the
Connes--Consani--Moscovici (CCM) framework, the only candidate that survives the structural
collapse.

\subsection{The governing thesis}
The opening survey (Phase 1) established the organizing thesis: \emph{every analytic front on
RH reduces to a positivity statement} --- Weil's explicit-formula functional $\ge0$,
equivalently Li's $\lambda_n\ge0$, Bombieri--Lagarias, the reality of the zeros. Any honest
attack must either (a) prove a positivity no one can prove ($=$ RH), or (b) be circular, or (c)
live in the wrong region. The program's value is in classifying which of (a)--(c) each attempt
falls into.

\subsection{Arc A --- the $\omega$-class / resonance program (Phases 1--6)}
Arc A attacked the Lindel\"of side via the $\omega$-class decomposition (organizing $n$ by
$\omega(n)$).
\begin{itemize}[leftmargin=*]
\item \textbf{P2.} $B(N)\le N^\varepsilon\Rightarrow$ Lindel\"of. \NOGO{} (NG-A1): the target is
stronger than Lindel\"of and requires it as a hypothesis --- circular.
\item \textbf{P3.} No power-law growth in the resonance sums; an unconditional onset bound
$N^*\gtrsim10^{14}$. A genuine, if negative, quantitative result.
\item \textbf{P4.} The \emph{coefficient architecture}, not the zero locations, is the
discriminant between $\zeta$ and its RH-violating cousins; a parity theorem for the Liouville
function organizes the even/odd $\omega$-classes. \PROVED{} (structural).
\item \textbf{P5--P6.} A putative ``phase transition'' at $N_c\approx1.18\cdot10^5$ was
\NOGO{} (NG-A6) --- refuted on canonical re-verification, honestly reframed. The lesson: the
resonance statistic $r$ is conditioning-dependent (mean zero), so any unconditional
cross-correlation claim dies by Montgomery--Vaughan (NG-A3).
\end{itemize}
Arc A also ruled out, with reasons, the parity-problem sieve (NG-A4), TDA as a detector
(NG-A7), the one-sided Gram artifact $\lambda_{\min}\approx0.866$ (NG-A8), and the
criteria-as-paths family Li / Jensen--GORZ / B\'aez-Duarte (NG-A5) --- reformulations of RH,
not routes to it.

\subsection{Arc B --- the inverse operator and the rigorous Weil detector (Phase 7)}
Arc B produced the program's strongest \emph{unconditional} results. The object is the
localized Weil quadratic form $Q=M_{\rm zeros}-M_{\rm arith}$ in a Hermite--Gauss basis at
height $T_0$, width $\sigma$, dimension $J$.
\begin{itemize}[leftmargin=*]
\item \textbf{Front A \PROVED.} An off-line zero forces $\lambda_{\min}(Q)<0$ with explicit
$\delta^2$ constants ($\delta=$ distance from the critical line). The form sees RH violations.
\item \textbf{Front B \PROVED{} (unconditional, PNT only).} The truncation floor obeys
$\eta(X)\le A\,e^{-\alpha(\log X)^2}$: the finite-$X$ error is super-polynomially small, no RH.
\item \textbf{Front C.} The detector is technical; the \emph{uniform} passage $Q_X\to Q_\infty$
is RH-equivalent --- the first clean appearance of the master wall.
\end{itemize}
The compression question (0A2) --- is $Q=P_J\mathcal T P_J$ a faithful compression of the Weil
operator $\mathcal T$? --- found both scenarios supported, leaving the named wall
$(LB):\ \inf\operatorname{spec}(\mathcal T)\ge0\iff\lambda_{\min}(Q_\infty)\ge0\ \forall
f=$ Weil positivity $=$ RH (if 0A2 holds). \WALL

\subsection{The $b_j$ lower-bound program and local--global (Phases 22--28)}
\begin{itemize}[leftmargin=*]
\item \textbf{Phases 22--23 (Paley--Wiener barrier).} Arithmetic information from
$\operatorname{Re}s>1$ provably does not propagate into $(1/2,1)$. This is the central wall
\textbf{MW-2}.
\item \textbf{Phases 24--25 (fifteen-mechanism audit).} Fifteen classical routes to a
\emph{lower} bound on $b_j$, each refuted with its failure mode \NOGO{} (NG-B1\dots B8 and more):
Korobov--Vinogradov and de Bruijn--Newman give the wrong direction; Deuring--Heilbronn lives
near $\operatorname{Re}s=1$; Beurling--Nyman, Tur\'an-for-$\xi_0$, Li/Krein--LP are circular;
Baker is the wrong category; Montgomery pair correlation and the large sieve carry no off-line
signal. Universal cause: MW-2.
\item \textbf{Phase 26.} Explicit eigenvectors $x^{b_j+i\gamma_j}\notin L^2(C_{\mathbb Q})$
\NOGO{} (NG-C3): the norm diverges; they are distributions.
\item \textbf{Phase 27 (local--global).} $Q=Q_\infty\prod_pQ_p$ fails on two walls \NOGO{}
(NG-C1): the zero side does not factor over primes (no Frobenius-at-$p$ for $\zeta/\mathbb Q$,
\textbf{MW-3}); Hasse--Minkowski fails for $\infty$-dim Hermitian forms (\textbf{MW-5}). Its
obstruction equals Phase 26's (NG-C2).
\end{itemize}

\subsection{The broad exploration (Phases 32--59)}
Phases 32--59 swept the structural landscape: CCM absolute geometry (32), de Bruijn--Newman
(33), five fronts and ABC forms (35--36), a physics/dynamics import (37, 53), the $G_1/G_2$
interface (39--41), Hodge-theoretic and foliated $\operatorname{Spec}\mathbb Z$ dynamics
(42--43), creative rupture / new mathematics (44), a Fredholm route (45), ``crossing the wall''
(49), a Diophantine thread (50), functor and homotopy-inertia constructions (51--52), index
dynamics (54), the ``two arrows'' / ``two towers'' (55--56), and the total audits and closure
(46--48, 58--59). The honest verdict: no genuinely new structure escaped the collapse --- which
made the pivot to CCM forced, not fashionable.

\subsection{The seven-fold wall and the meta-search collapse theorem}
A single obstruction reappears under seven names: \textbf{MW-1} Weil positivity; \textbf{MW-2}
arithmetic propagation ($\operatorname{Re}s>1\not\to(1/2,1)$, central); \textbf{MW-3}
Hasse--Minkowski in $\infty$-dim; \textbf{MW-4} the Hilbert--P\'olya \emph{wrong-sign} wall ---
every unconditional positivity tool gives a lower bound, RH needs an upper; \textbf{MW-5}
Connes--Consani arithmetic-site; \textbf{MW-6} uniform spectral gap, itself $=$ RH. An
adversarial novel-structure search, literature-merged and attacked to destruction, proved the
collapse theorem: \emph{every RH-directed structure reduces to CAP (positivity), SURF
(Connes--Consani geometry, RH-equivalent), or F (statistics)}. The literature reproduces it ---
Kurasov--Sarnak Fourier quasicrystals (reality of support $=$ Lee--Yang $=$ stable polynomial
$=$ Hermite positivity $=$ CAP), Deninger ($=$ CAP), Nyman--Beurling/Li ($=$ C), de
Branges/Sonine ($=$ Hermite--Biehler reality), Berry--Keating $xp$ ($=$ Hilbert--P\'olya),
CFKRS/RMT ($=$ F). Two escape candidates destroyed: \textbf{(A)} crystalline rigidity $=$
Lee--Yang $=$ positivity; \textbf{(B)} functional-equation edge-anchor $=$ the Weil/de Branges
form $=$ CAP. The center is reached by definiteness, never by an anchor.

\subsection{The forced pivot: CCM spectral triples (Phase 60)}
The collapse leaves one live, non-circular candidate. F cannot decide reality; CAP is the wall.
SURF --- the CCM spectral triples --- is distinguished by \textbf{MW-4 inverted}: the operators
are self-adjoint \emph{by construction}, so the spectrum is real automatically and the open
problem is operator convergence $H_x\to H_\infty$, \emph{not} the sign. This converts the
wrong-sign wall into an approximation question --- the one place the master obstruction does not
obviously apply. Phase 60 implements CCM faithfully and pushes the convergence/positivity
question as far as it goes, yielding the reduction chain of \S3.

\section{Paths explored and no-go results}

\subsection{Arc A ($\omega$-class / resonance)}
\begin{longtable}{l p{5.4cm} p{6.4cm}}
\toprule \textbf{Tag} & \textbf{Attempt} & \textbf{Why it fails} \\ \midrule \endhead
NG-A1 & $B(N)\le N^\varepsilon\Rightarrow$ Lindel\"of & Target stronger than Lindel\"of; circular. \\
NG-A2 & Extrapolate $\alpha_B\to0$ & $\alpha_B\approx0.42$ vs threshold $<0.31$; never crosses. \\
NG-A3 & Unconditional cross-covariance as signal & Montgomery--Vaughan: $\approx0$ unconditionally. \\
NG-A4 & Sieve to separate $\omega$-parity & Parity problem blocks all standard sieves. \\
NG-A5 & Li / Jensen--GORZ / B\'aez-Duarte as a path & Reformulations of RH; circular. \\
NG-A6 & Phase transition at $N_c$ & Refuted on re-verification; no transition. \\
NG-A7 & TDA as a zero detector & $H_0$ is $\delta^2$-insensitive; $H_1$ no signal. \\
NG-A8 & $\lambda_{\min}\approx0.866$ as positivity & One-sided Gram trivially positive. \\
\bottomrule
\end{longtable}

\subsection{The $b_j$ lower-bound mechanisms (Phases 24--25)}
\begin{longtable}{l p{5.4cm} p{6.4cm}}
\toprule \textbf{Tag} & \textbf{Attempt} & \textbf{Why it fails} \\ \midrule \endhead
NG-B1 & Fifteen classical mechanisms en bloc & MW-2: $\operatorname{Re}s>1$ info does not reach $(1/2,1)$. \\
NG-B2 & Korobov--Vinogradov zero-free region & Gives $b_j\le1/2-c/(\log\gamma_j)^{2/3}$ --- upper bound. \\
NG-B3 & Deuring--Heilbronn repulsion & Operates near $\operatorname{Re}s=1$; no anti-Siegel off the line. \\
NG-B4 & Beurling--Nyman distance & $d>F(b_j)$ from below $\iff$ bounding $b_j$ --- circular. \\
NG-B5 & de Bruijn--Newman $\Lambda\ge b_j^2/2$ & Inverted: gives $b_j\le\sqrt{2\Lambda}$, trivial. \\
NG-B6 & Tur\'an inequalities for $\xi_0$ & Satisfied for all $b_j>0$ --- no constraint. \\
NG-B7 & Baker (linear forms in logs) & Applies to algebraic numbers; $\gamma_j$ transcendental. \\
NG-B8 & Montgomery pair correlation & On-line spacings only; no off-line signal. \\
\bottomrule
\end{longtable}

\subsection{Local--global and distributional (Phases 26--27)}
\begin{longtable}{l p{5.4cm} p{6.4cm}}
\toprule \textbf{Tag} & \textbf{Attempt} & \textbf{Why it fails} \\ \midrule \endhead
NG-C1 & $Q=Q_\infty\prod_pQ_p$ (Hasse--Minkowski) & Zero side does not factor over primes (MW-3); H--M fails $\infty$-dim (MW-5). \\
NG-C2 & Phase 27 as alternative to 26 & Same obstruction, two languages. \\
NG-C3 & Eigenvectors $x^{b_j+i\gamma_j}\in L^2$ & Norm diverges; distributions, $Q(f_j,f_j)=\infty$. \\
\bottomrule
\end{longtable}

\subsection{The structural walls and the meta-search}
$MW$-1 Weil positivity; $MW$-2 arithmetic propagation (central); $MW$-3 Hasse--Minkowski
$\infty$-dim; $MW$-4 Hilbert--P\'olya wrong sign; $MW$-5 Connes--Consani arithmetic site; $MW$-6
uniform spectral gap ($=$ RH) --- the same wall from seven directions. Every RH-directed
structure (corpus or literature) collapses to CAP (positivity), SURF (Connes--Consani,
RH-equivalent), or F (statistics); the two purpose-built escapes (crystalline rigidity $=$
Lee--Yang; functional-equation anchor $=$ Weil/de Branges) are destroyed.

\subsection{Phase-60 no-gos (the CCM scaffolding)}
\begin{longtable}{p{6.3cm} p{8cm}}
\toprule \textbf{Path} & \textbf{Why it fails} \\ \midrule \endhead
$H1$ proves $QW$ PSD via Perron--Frobenius & PF gives ground-state structure, not positive-definiteness. PSD $=$ RH. \\
Limit operator constant $c=-\psi''(1/4)/8$ & Archimedean curvature subdominant by $\sim\lambda\log^2\lambda$; refuted by $(2\lambda)^{-1}D_\lambda(s/\log\lambda)\to\cos2s$. \\
Hankel edge model, $\gamma\to1$, scale $2\lambda$ & False premise $|\ve_0|\sim2\lambda$; in fact $\sim2.4\cdot10^{-48}$ at $\lambda^2{=}11$. \\
Dirichlet box / standard Agmon at $e^{-L}$ & Box gives $1/\log^2\lambda$, tunneling $\ve_1/\ve_0\to\infty$; neither $e^{-L}(k{+}1)^2$. \\
Magnitude methods for $\operatorname{sign}\ve_0$ & Marginality $\sqrt2/\pi$; depends only on $\{|\Lambda(n)|\}$; DH same magnitudes, $\ve_0<0$. \\
Frequency law $T_p=\pi/\log p$ for $\delta a_n$ & Confirmed only for $p=2$; fits reject $2\log p$ for $p\ge3$. \\
Project $J_S$ out for Sonin & $\xi_S$ cyclic $\Rightarrow$ complement $\{0\}$; Sonin in position space. \\
$E_S=1/\Lambda_S$ as de Branges function & Correct modulus, not Hermite--Biehler; needs the functional equation. \\
\bottomrule
\end{longtable}

\noindent\textbf{Convergent lesson.} Every route to the zeros --- $\omega$-resonance, the
localized detector, $b_j$ bounds, local--global, the CCM prolate, the de Branges/Sonin side ---
terminates at the same place: the reality of the $\gamma_n$, equivalently Weil positivity,
equivalently RH. The sign is a \emph{phase} phenomenon, immune to magnitude estimates, to
local-operator models, and to any structure that does not encode the zeros by construction.

\section{The proof chain}

\subsection{Architecture of the proof: where RH would be proved}\label{sec:arch}

Before the technical steps, we explain the shape of the argument as a whole, so that the role
of each step --- why it is there, what it receives from the step before, and what it hands to
the step after --- is visible in advance. The reader who keeps the following six-sentence
skeleton in mind will not get lost in the lemmas.

\emph{The strategy is to convert RH from a statement about the location of zeros into a
statement about the sign of a single number, and then to prove every link except that sign.}
The object that carries the conversion is a one-parameter family of quadratic forms
$A_\lambda$, $\lambda>1$, built from the Weil explicit formula and localized to band-limited
functions of bandwidth $\lambda$. Its construction and meaning are \textbf{Step 1}: $A_\lambda$
is simultaneously (a) a concrete convolution operator with an explicit symbol $a(t)=\Om(t)-
D_\lambda(t)$, and (b), by the Weil explicit formula, a sum over the zeros $\rho$ of $\zeta$ ---
so its positivity is \emph{equivalent} to the zeros being real. This double identity is the
hinge of everything: the analytic side (the symbol) is computable, the arithmetic side (the
zeros) is what we want to control.

The chain then proceeds as follows. \textbf{Step 2} analyses the bottom of the spectrum of
$A_\lambda$: it shows the lowest eigenvalue $\ve_0(\lambda)$ is simple with a clean
eigenfunction, and --- this is the new and delicate part --- that the \emph{gap} above it does
not collapse as $\lambda\to\infty$, so that the family $A_\lambda$ has a stable limiting shape.
\textbf{Step 3} proves that, because $A_\lambda$ is a compression of a convolution
(difference-kernel) operator, the entire function $\hat\xi_\lambda$ attached to its ground state
has only real zeros at each finite $\lambda$. \textbf{Steps 4--5} are the reduction proper: the
finite-$\lambda$ object approximates the completed zeta function $\Xi$ with explicit error, and
real-rootedness passes to the limit \emph{if and only if} a single localized quantity
$\ve_\loc(\lambda)$ decays faster than $\lambda^{-1/2}$. At this point RH has been turned into
an inequality about one number. \textbf{Steps 6--7} certify that the reduction is not circular
and not vacuous: a known RH-false model (Davenport--Heilbronn) is correctly \emph{rejected} by
the same machinery, and the classical Hurwitz/Laguerre--P\'olya theory supplies the limiting
logic. \textbf{Steps 8--10} identify the limit operator inside the CCM prolate/zeta-spectral-
triple framework and reconstruct, as exact identities, the arithmetic content (periodic-orbit /
Gutzwiller sums over primes) that any correct operator must carry; these steps are RH-neutral
scaffolding --- they make the operator concrete without assuming the answer.

\textbf{This is where RH would be proved.} After Step 5, RH is equivalent to
$\ve_\loc(\lambda)=o(\lambda^{-1/2})$, equivalently (Step 12.1) to the positivity
\[
(P):\qquad \ve_0(\lambda)\ge 0\quad\text{for all }\lambda>1 .
\]
Every step 1--10 is in place to make $(P)$ a faithful, non-circular, phase-sensitive
restatement of RH and to surround it with an explicit operator. \textbf{Steps 11--12 are the
wall}: they show that $(P)$ is exactly as hard as RH itself --- it is reached identically from
the primes, from the zeros, from the Sonin space of the prolate operator, and from de Branges
canonical-system positivity, and the marginality theorem (Step 12.2) proves that no
\emph{magnitude} method can decide its sign, because the sign is a \emph{phase} invariant. A
proof of RH along this route is therefore precisely a proof of $(P)$: one must show the
localized Weil form never goes negative, and the program's contribution is to have reduced RH
to that one statement with everything around it proved, and to have identified --- in
Lemma~\ref{lem:23F} --- the single new analytic fact (the arithmetic deformation of the
spectral ladder) on which the \emph{stability} of the whole picture depends.

Throughout, a status mark records the epistemic status of each claim (see the abstract), so the
boundary between what is proved and what is the wall is never blurred.

\paragraph{Notation.} $\lambda>1$, $L=2\log\lambda$, $\Omega=\log\lambda$; $\Lambda$ von
Mangoldt, $\psi(X)=\sum_{n\le X}\Lambda(n)$; $\gamma_\rho$ the ordinates of the nontrivial
zeros, $T^*=2\pi\lambda^2$; $\Om(t)=-\log\pi+\operatorname{Re}\psi_\Gamma(\tfrac14+\tfrac{it}2)$
the archimedean symbol, $D_\lambda(t)=\sum_{n\le\lambda^2}\Lambda(n)n^{-1/2}\cos(t\log n)$ the
prime symbol; $E_\lambda$ the band-limited functions of exponential type $\Omega$, spectral
support in $[-L/2,L/2]$.

\subsection{Step 1 --- The Weil quadratic form and its symbol \PROVED}

\emph{Purpose, and why it comes first.} Everything downstream is an analysis of one object, so
we must first build it and establish the one property that gives it meaning: that its positivity
is equivalent to RH. The Weil explicit formula expresses a sum over the zeros of $\zeta$ as a
sum of an archimedean term and a sum over prime powers. Read as a quadratic form in a test
function $f$, it becomes a Hermitian form $A_\lambda$ whose value is, on one hand, a transparent
analytic expression and, on the other hand, literally $\sum_\rho|\hat f(\gamma_\rho)|^2$ when
the zeros are real. We localize to $f\in E_\lambda$ (band-limited, bandwidth $\lambda$) so that
only finitely many primes ($n\le\lambda^2$) enter and the form is a genuine compact-interval
operator. Concretely, for $f\in E_\lambda$,
\[
A_\lambda(f,f)=\int\Om(t)\,|\hat f(t)|^2\,dt-\int D_\lambda(t)\,|\hat f(t)|^2\,dt ,
\]
the difference of the archimedean symbol $\Om$ and the prime symbol $D_\lambda$ tested against
the power spectrum $|\hat f|^2$. The three lemmas below establish, in turn, the three faces of
$A_\lambda$ we will use later: its kernel (so that Step~3's convolution theorem applies), its
Fourier symbol (so that Steps~2 and~12 can do spectral perturbation theory on $a(t)$), and its
Weil identity (the RH-equivalence that justifies the entire enterprise).

\begin{lemma}[Schwartz kernel]
$A_\lambda$ is the compression to $[-L/2,L/2]$ of convolution by $D_\zeta(|x-y|)$, where
$D_\zeta(u)=w_{\mathrm{arch}}(u)-\sum_{n\le\lambda^2}\Lambda(n)n^{-1/2}\delta(u-\log n)$.
\end{lemma}
\begin{proof}
From $Q(f,g)=\int_{-L}^{L}(g\star f)(y)D_\zeta(|y|)dy$, $(g\star f)(y)=\int\overline{g(x)}f(x+y)dx$,
$h(u):=D_\zeta(|u|)$ even, change variables $(x,y)\mapsto(x,y'{=}x{+}y)$:
$Q(f,g)=\iint\overline{g(x)}f(y')h(y'-x)dx\,dy'=\int\overline{g(x)}\big(\int k(x,y')f(y')dy'\big)dx$,
$k(x,y')=D_\zeta(|x-y'|)$. Restriction gives the compression.
\end{proof}
\begin{lemma}[Fourier multiplier]
$\hat h(t)=2\int_0^\infty D_\zeta(u)\cos(tu)du=\Om(t)-D_\lambda(t)=:a(t)$.
\end{lemma}
\begin{proof}
The prime part gives $D_\lambda(t)$; the archimedean part
$2\int_0^\infty w_{\mathrm{arch}}(u)\cos(tu)du=\Om(t)$ by the explicit-formula relation.
\end{proof}
\begin{lemma}[Weil identity]\label{weil}
$A_\lambda(f,f)=\sum_\rho\hat f(\gamma_\rho)\overline{\hat f(\overline{\gamma_\rho})}$; under RH,
$=\sum_\rho|\hat f(\gamma_\rho)|^2\ge0$.
\end{lemma}
\begin{proof}
Weil's explicit formula on $g\star f$, support $2\Omega=L$ truncating primes to $n\le\lambda^2$;
archimedean and pole terms give $w_{\mathrm{arch}},\mathcal P$; the zero sum is indefinite absent
reality of $\gamma_\rho$.
\end{proof}
Finally we record the form in \emph{carr\'e-du-champ} (energy) coordinates, which is the
representation Step~2 needs: symmetrizing the kernel splits $A_\lambda$ into a mass term and two
non-negative Dirichlet (``energy'') forms,
\[
A_\lambda(f,f)=\underbrace{(\Om(0)-P_\lambda)\|f\|^2}_{\text{mass}}
+\underbrace{\mathcal E_\theta(f)}_{\text{archimedean jumps}}
+\tfrac12\underbrace{\mathcal E_{\rm prime}(f)}_{\text{prime jumps}},\qquad
\mathcal E_\theta,\mathcal E_{\rm prime}\ge0 .
\]
The point of this rewriting is that it separates what can make $A_\lambda$ negative (the scalar
mass coefficient $\Om(0)-P_\lambda$) from what is manifestly positive (the two energy forms).
Step~2 exploits exactly this separation. \PROVED

\subsection{Step 2 --- Structure of the ground state \PROVED{} ($\lambda$ finite)}

\emph{What Step~2 must deliver, and why.} For the reduction in Step~5 to pass to the limit we
need two things about the bottom of the spectrum of $A_\lambda$: that the ground state is
\emph{clean} (simple, isolated, with a sign-definite eigenfunction, so that the entire function
$\hat\xi_\lambda$ of Step~3 is well defined) and that it is \emph{stable} (the spectral gap above
$\ve_0$ does not close as $\lambda\to\infty$, so that the limit $\lambda\to\infty$ in Step~5 does
not degenerate). The first is a soft Perron--Frobenius argument and is unconditional; it uses
only the positivity of the energy forms isolated in Step~1. The second is the genuinely hard,
genuinely new content of this paper, and occupies Lemmas~\ref{lem:23A}--\ref{lem:23F} and
Theorem~\ref{thm:23} below. We begin with the positivity structure that makes the semigroup of
$A_\lambda$ positivity-preserving.

\begin{lemma}[L\'evy--Khintchine positivity]
$M_\theta\succeq\Om(0)I$, and $M_\theta-\Om(0)I$ is a jump Dirichlet form; its semigroup is
positivity-preserving.
\end{lemma}
\begin{proof}
$\Om(t)-\Om(0)=\int_0^\infty(1-\cos t\xi)d\nu(\xi)$, $d\nu=2e^{-\xi/2}(1-e^{-2\xi})^{-1}d\xi\ge0$
(verified $\int\min(1,\xi^2)d\nu<\infty$). Hence
$\langle f,M_\theta f\rangle-\Om(0)\|f\|^2=\tfrac12\iint|f(w)-f(w+\xi)|^2d\nu\,dw\ge0$.
\end{proof}
\begin{theorem}\label{thm:22}
At each finite $\lambda$, $\ve_0(\lambda)$ is simple, isolated, with even constant-sign
eigenvector $\xi$.
\end{theorem}
\begin{proof}
$A_\lambda=M_\theta-V$, $V=\sum\Lambda(n)/\sqrt n\,T_{\log n}$. $e^{-tM_\theta}$ and (since
$\Lambda\ge0$) $e^{tV}$ positivity-preserving; Trotter $\Rightarrow$ Perron--Frobenius
$\Rightarrow$ bottom simple, isolated, constant sign; $\iota:x\mapsto L-x$ forces evenness.
\end{proof}
\begin{remark}[uniformity --- the rescaled low spectrum]
The uniformity $\liminf_\lambda(\ve_1-\ve_0)/|\ve_0|>0$ is not delivered by the
Perron--Frobenius argument, which is $\lambda$-by-$\lambda$. With the validated CCM engine
(dps $40$, primes included) the rescaled low spectrum of $A_\lambda$ obeys
$\ve_k/\ve_0\to(k+1)^2=4,9,16,25,\dots$, with $\ve_2/\ve_0\to9.0$ very cleanly and
$\ve_1/\ve_0\in[3.51,4.19]$ robustly across $\lambda\in[3,12]$, $N\in[8,24]$. This is the
\emph{positive} prolate near-radical spectrum (the Slepian $\chi_m$), not the negative
Sonin spectrum of Step 11. Two naive routes are refuted (NO-GO): (i) archimedean skeleton
$+$ bounded prime perturbation (primes off $\Rightarrow$ $M_\theta$ indefinite,
$\widetilde\Omega(0)\approx-5.37<0$, no ladder --- the ladder is constitutively
arithmetic); (ii) free Dirichlet box (spectrum $(k+1)^2$-like but eigenvectors are not box
modes). The correct identification follows.
\end{remark}

\emph{Strategy for the uniform gap.} We must show $\liminf_\lambda(\ve_1-\ve_0)/|\ve_0|>0$. The
obstacle is that the absolute scale of the low eigenvalues is super-exponentially small
($\rho_\lambda\asymp e^{-cL}$), so any argument that tries to track $\ve_0$ itself is hopeless;
we must work with \emph{ratios}, in which the unknown scale cancels. The plan has three moves.
First (Lemma~\ref{lem:23A}, the Doob identity) we record the exact change of variables that
turns the candidate limit operator into a Dirichlet Laplacian, so that ``the ratios are
$m^2$'' and ``the eigenfunctions are Chebyshev'' become the same statement. Second
(Lemma~\ref{lem:23B}, scale-free Parter) we prove the general principle that a Dirichlet-type
form has ratios $\to m^2$ \emph{independently of its overall scale} --- this is what licenses us
to ignore $\rho_\lambda$. Third (Lemma~\ref{lem:23F}, the CCM shorted Dirichlet quotient) we
state the one fact that is genuinely about the explicit arithmetic of the Weil form: that the
conditioned (shorted) form on the slow sector \emph{is} a Dirichlet form. Theorem~\ref{thm:23}
then combines the three into the uniform gap, and Corollary~\ref{cor:23G} feeds that gap into
Step~3 to make real-rootedness survive the limit. Lemma~\ref{lem:23F} is the only assumed
ingredient; the remark after it explains precisely why it is new and how far its proof has been
carried.

\begin{lemma}[Doob identity --- PROVED]\label{lem:23A}
For $f=\sin\theta\cdot v$, $\int_0^\pi|f'|^2-\int_0^\pi|f|^2=\int_0^\pi|v'|^2\sin^2\theta$.
\end{lemma}
\begin{proof}
$f'=\cos\theta\,v+\sin\theta\,v'$; the cross term
$\int2\sin\theta\cos\theta\,vv'=-\int\cos2\theta\,v^2=-\int v^2+2\int\sin^2\theta\,v^2$ (by
parts, boundary terms vanish); collecting with $\cos^2\theta-1=-\sin^2\theta$ everything
cancels to $\int\sin^2\theta\,v'^2$. Hence the Dirichlet Laplacian (eigenvalues $m^2$,
eigenfunctions $\sin m\theta$), shifted by its ground energy $1$ and transformed by
$s_1=\sin\theta$, becomes $L_\infty=-\sin^{-2}\theta\,\partial_\theta(\sin^2\theta\,\partial_\theta)$
(Chebyshev $U_k(\cos\theta)=\sin((k+1)\theta)/\sin\theta$, eigenvalues $k(k+2)$).
\end{proof}
\begin{lemma}[scale-free Parter --- PROVED]\label{lem:23B}
If $F_N=\tau_N D_D^*C_N D_D$ with $D_D$ the Dirichlet first-difference and $C_N\to c\,I$ in
low frequencies, then $\lambda_m(F_N)=\tau_N c\cdot4\sin^2(\pi m/2(N+1))(1+o(1))$, so
$\lambda_m(F_N)/\lambda_1(F_N)\to m^2$ for fixed $m$ --- independent of the scale $\tau_N$
(it cancels in the ratio). This is why the super-exponential $\rho_\lambda\asymp e^{-cL}$
need never be computed.
\end{lemma}
\begin{lemma}[CCM shorted Dirichlet quotient --- ASSUMED; the genuine residual]\label{lem:23F}
In the CCM objects: let $\mathcal{E}(f)(x)=x^{1/2}\sum_{n>0}f(nx)$ be the conditioning map
(Connes--Consani, arXiv:2106.01715, eq.~3.1), let $\varphi_n$ be the prolate combinations
vanishing at $0$, and let $\varepsilon_n$ be the Gram--Schmidt orthonormalization of
$\{\mathcal{E}(\varphi_n)\}$ (eq.~3.4), spanning the prolate projection $\Pi(\lambda,k)$
(Def.~3.1). Then for every fixed $N$,
$\rho_\lambda^{-1}\langle\varepsilon_m,QW_\lambda\varepsilon_n\rangle\to m^2\delta_{mn}$
$(1\le m,n\le N)$, i.e. the shorted/quotient form of $QW_\lambda$ on the near-radical is
the Dirichlet form $-\partial_\theta^2$. This is not a soft consequence of
Lemmas~\ref{lem:23A}/\ref{lem:23B}; it is a genuine theorem about the explicit
$\mathcal{E}$, here assumed. Validated: (i) eigenvectors $\varepsilon_0,\varepsilon_1,\varepsilon_2$
are the Dirichlet sine modes (arXiv:2106.01715, Figs.~22--24), matching the true
$QW_\lambda$ eigenfunctions (Figs.~26--37); (ii) Rayleigh quotients
$R(U_k)\to k(k+2)=3,8,15$ (exact engine, $\lambda=7,10,13$); (iii) overlaps
$0.99,0.80,0.51$; (iv) $\ve_k/\ve_0\to(k+1)^2$ robust in $\lambda,N$.

\emph{Status (2026-06-21; head-to-head vs the CCM prolate Jacobi --- Connes-validated, his
version pending).} The ladder is a property of the Weil form $QW_\lambda$, \emph{not} of the
bare prolate $W_\lambda$, and is not in the CCM prolate papers. The exact CCM Jacobi
(arXiv:2310.18423, Prop.~3.5: $b_n=-8n^2-2n-\tfrac34+2\pi\lambda^2(8n+1)$) gives, under Schur
complement, the \emph{linear} ladder $8n+1$ (normalized $0,1,2,3,4$); the engine $QW_\lambda$
gives the \emph{quadratic} ladder $\to(k+1)^2$. So the quadratic ladder is constitutively
arithmetic, living in $QW_\lambda-W_\lambda$ (prime $+$ archimedean terms): \textbf{2.3.F is a
new theorem, not a citation.} Route (A) diagnosis (partial, not closed): in the symbol
$a(t)=\widetilde\Omega(t)-D_\lambda(t)$ the curvature $a''(0)=\widetilde\Omega''(0)+M_2$ with
$M_2=\sum\Lambda(n)(\log n)^2/\sqrt n\sim8\lambda\log^2\lambda$ dominates, so $\tfrac12 M_2 t^2
\leftrightarrow-\partial_w^2$ is the local Laplacian with arithmetic coefficient. Open piece:
the marginality $\ve_0\approx0$ does not follow from the bare symbol ($a(0)\approx-17<0$); it
comes from the regularization $WR$/near-radical, which also supplies the Dirichlet
confinement. Assumed pending Connes' version.
\end{lemma}

\begin{remark}[Phase 2.3.F-CLOSE (2026-06-23): reduction to one sign-free residual + eight refuted shortcuts]\label{rem:23Fclose}
A dedicated campaign reduced Lemma~\ref{lem:23F} to a single sharply-stated residual and mapped the obstruction exhaustively. Throughout, the operator is the \emph{exact} engine $QW_\lambda$ (never the symbol $a(D)$, which is the wrong, indefinite operator).

\emph{\textbf{(R1) Reduction (PROVED assembly).}} Let $u_0>0$ be the ground state (Theorem~\ref{thm:22}, Perron--Frobenius, unconditional at finite $\lambda$). Doob-conjugation by $u_0$ gives the quotient form $G_\lambda(v,v)=A_\lambda(u_0v,u_0v)-\ve_0\|u_0v\|^2$, self-adjoint for the weight $w=u_0^2$, with eigenpairs $(\ve_k-\ve_0,\;v_k=f_k/u_0)$, \emph{manifestly positive} (a carr\'e-du-champ; the sign is automatic \emph{after} conjugation by the proved positive ground state). Define
\[\textbf{Lemma 2.3.F-Loc:}\quad G_\lambda v=-w^{-1}(C_\lambda\,w\,v')'\ \text{in the bulk, with}\ C_\lambda(y)/\overline{C_\lambda}\to1\ \text{as}\ \lambda\to\infty.\]
Then $C_\lambda\to\mathrm{const}$ together with Parter (Lemma~\ref{lem:23B}, scale-free, cancels $\rho_\lambda\asymp e^{-cL}$), Dirichlet-BC (Lemma~2.3.F-Box, time-limiting forces $H^1_0$), and Doob (Lemma~\ref{lem:23A}) yield $(\ve_k-\ve_0)/|\ve_0|\to(k+1)^2-1$. \emph{Every ingredient except 2.3.F-Loc is proved}; 2.3.F-Loc is sign-free (it constrains the shape of $C_\lambda$, never $\operatorname{sign}\ve_0$).

\emph{\textbf{(R2) Validation (E70).}} $C_\lambda(y)$ is flat in the bulk with $\mathrm{std}/|\mathrm{mean}|=0.222\to0.146\to0.137$ \emph{decreasing} in $\lambda=7,9,11$; the local Rayleigh $\overline{C_\lambda}\int w v_k'^2/\int w v_k^2$ reproduces the true Doob eigenvalues ($R/e=1.03$--$1.10$), ratios $[1,2.69,5.18]$ vs.\ $k(k+2)$ normalization $[1,2.667,5]$; the balance $\mathrm{ARCH}=\mathrm{PRIME}$ holds to $>16$ digits. \emph{Davenport--Heilbronn falsifier}: $C(y)$ non-flat ($\mathrm{std}/|\mathrm{mean}|=33$), local Rayleigh fails ($[1,1.6,2.28]$), $\mathrm{ARCH}\ne\mathrm{PRIME}$ ($\ve_0<0$).

\emph{\textbf{(R3--R8) Eight refuted shortcuts.}} Every approximate/perturbative/local route to 2.3.F-Loc fails, each for the same root reason --- the ladder lives in the $e^{-cL}$ residual of the \emph{exactly-balanced} Weil form, below any approximation:
\begin{enumerate}\setlength\itemsep{0pt}
\item[(R3)] \emph{Raw second moment} (E71): $c_2=\sum\Lambda(n)n^{-1/2}(\log n)^2$ is large-jump dominated (small primes give $1\%$); non-local, predicts the \emph{opposite} (centre-deficit) profile.
\item[(R4)] \emph{Autocorrelation taper} $(1-s/L)^p$ (E72): recovers the magnitude of the flatness but not its $\lambda$-trend; necessary, not sufficient.
\item[(R5)] \emph{Doob-tempered moment} (E74): with the correct tempering $u_0(x\pm\log n)/u_0(x)$ the pointwise multiplier is edge-divergent (centre/edge $=0.08$, worsening) --- third failure of the local-coefficient class; the ladder is \emph{spectral}, not pointwise.
\item[(R6)] \emph{Bare-prolate deformation} (E75a): the Connes Jacobi (Prop.~3.5) gives linear $0,1,2,3,4$; interpolating $\mathrm{ARCH}-t\,\mathrm{PRIME}$ leaves the ladder \emph{frozen collapsed} for all $t\in[0,1)$, the quadratic ladder appearing only at $t=1$ exactly.
\item[(R7)] \emph{Near-marginal limit} (E75b): at $t=0.9999$ ($\ve_0=-4\!\cdot\!10^{-4}$, minute) the ladder is still collapsed; the quadratic ladder is a knife-edge at $\ve_0=0$ to precision better than $e^{-cL}$. \emph{The marginality theorem (12.2) does not by itself force the ladder.}
\item[(R8a)] \emph{Jacobi unmasking} (E77): removing the scaling term $2\pi\lambda^2(8n+1)$ to expose the pre-existing $-8n^2$ gives \emph{sub-linear} $0,1,1.76,2.39$, not $(k+1)^2$; the arithmetic does not merely cancel the scaling term.
\item[(R8b)] \emph{Classical-prolate eigenfunctions} (E78): $\mathrm{overlap}(f_k^{QW},\psi_k^{\text{sinc-prolate}})=0.04$--$0.5$; the $QW$ eigenfunctions are prolate-\emph{concentrated} but \emph{not} the classical prolate-spheroidals, so classical prolate WKB cannot be borrowed for the carrier affinity.
\end{enumerate}

\emph{\textbf{The sharp crux.}} From E22 the near-radical eigenfunctions are pure Dirichlet sines $\sin((k+1)\theta)$ in the carrier coordinate $\theta$, with $u_0\propto\sin\theta$ (so $f_k/f_0=U_k(\cos\theta)$). Hence \textbf{2.3.F-Loc $\iff$ ``$\theta(y)$ is asymptotically affine as $\lambda\to\infty$''} ($\iff C_\lambda\to\mathrm{const}$ $\iff$ the near-radical operator is constant-coefficient $-\partial_\theta^2$). The carrier affinity residual is measured at $6.5\%\to3.8\%$, decreasing. There is \emph{no} exact algebraic ladder (ratios $4.07,3.998,3.904$ at $\lambda=7,9,11$ are not exactly $4$): the statement is genuinely \emph{asymptotic}, requiring the prolate/Sonin analysis of the $QW$ deformation ($8n+1\to(k+1)^2$, Connes' route A) of the exactly-balanced residual --- the same machinery as, but in \emph{shape} (RH-neutral) rather than \emph{sign} (RH) form, the wall. This is the precise open task. \emph{(Net status: 2.3.F reduced to 2.3.F-Loc, validated, DH-falsified, eight shortcuts conclusively refuted, crux named; the residual prolate/Sonin asymptotics remain genuinely open.)}

\smallskip
\emph{\textbf{(N1--N3) Phase 2.3.F-CLOSE-II (2026-06-25): three new intrinsic/robust fronts, $\lambda=7\!\dots\!21$, DH-falsified throughout (engine\_cache.py; full log R9\_three\_new\_fronts\_RESULTS.md).}} Stepping away from per-eigenvalue attacks (fragile at scale $e^{-cL}$): \textbf{(N1) Intrinsic Jacobi via exact Lanczos on the Weil matrix $A$ (E79)} --- a both-parities cyclic Lanczos (mpmath, full reorthogonalization) \emph{reproduces the $(k+1)^2$ ladder exactly inside the residual} ($\lambda=9$: $[1,3.998,9.03,16.06,25.2,36.4]$, $\mathrm{rep\_err}=0$), \emph{defeating the obstruction of E76} (an internal orthogonal reduction is stable where an external sine basis injected $O(1)$ error). Measured structural insight: the intrinsic Jacobi has \emph{bounded $O(1)$ coefficients} ($\alpha_n\approx[0.79,3.69,0.85,3.68,\dots]$, parity-interleaved), \emph{not} the scale-dominated $b_n\approx2\pi\lambda^2\,8n$ of the bare prolate Jacobi (Prop.~3.5): the Weil+Doob balance \emph{strips the $2\pi\lambda^2$ scale}, leaving a bounded-coefficient Jacobi whose low spectrum is $(k+1)^2$ --- a clean reframing of route A as \emph{``bounded intrinsic Jacobi $\to$ Dirichlet $(k+1)^2$''}. \textbf{(N2) Spectral zeta / heat trace (E80)} --- with $\mu_k=\ve_k/\ve_0$, the invariants $Z(s)=\sum\mu_k^{-s}$, $\Theta(t)=\sum e^{-\mu_k t}$ match the Dirichlet $(k+1)^2$ invariants to $<0.5\%$ aggregate (far more stable than individual eigenvalues), DH failing catastrophically ($Z(s)$ error $4.3,10.8,19.9$); \emph{but $\mu_1$ oscillates around $4$ ($4.07,4.00,3.90,4.04,3.93,4.03,3.98,3.93$), not a power law}. \textbf{(N3) Liouville--Green carrier (E81)} --- small affinity residual ($0.5$--$2.5\%$) but \emph{flat in $\lambda$} ($1/L$ rate $\approx0$). \textbf{Envelope to $\lambda=21$ (E82): the deviation from $(k+1)^2$ does \emph{not} decay cleanly; it oscillates in a stable band} ($\mathrm{ladder\_err}=0.025,0.009,0.024,0.013,0.017,0.010,0.006,0.018$). \emph{\textbf{Verdict R9 (honest):}} the three fronts do \emph{not} close 2.3.F, but yield real progress --- a \emph{stable} computational route to the residual ladder (Lanczos), a \emph{sharper crux} (bounded intrinsic Jacobi $\to$ Dirichlet, $2\pi\lambda^2$ scale provably removed), \emph{triple-robust} confirmation of 2.3.F-Loc with clean DH falsification, and the obstruction \emph{precisely identified}: an \emph{irreducible $\sim\!1$--$2\%$ arithmetic fluctuation oscillating in $\lambda$} (prime-oscillation echo in the band $L=2\log\lambda$), so the ladder is exact \emph{in the mean} and the remaining content is a \emph{Ces\`aro/averaged} prolate--Sonin asymptotic.

\smallskip
\emph{\textbf{(R10) The closing route (2026-06-25): band edge of the bounded intrinsic Jacobi --- 2.3.F-Loc is RH-NEUTRAL, not the wall (E85; log R10\_bandedge\_closure.md).}} \emph{De-hedging correction: 2.3.F-Loc is already RH-neutral by R1.} After Doob conjugation by the positive ground state $u_0$ (Perron--Frobenius, unconditional at finite $\lambda$), $G_\lambda=A_\lambda(u_0\cdot,u_0\cdot)-\ve_0\|u_0\cdot\|^2\ge0$ is \emph{manifestly positive at every $\lambda$, regardless of RH}; the ratios $e_k/e_1\to k(k+2)$ are an \emph{unconditional} spectral question about $G_\lambda$, while the wall (\S12.1) is the \emph{sign} of $\ve_0$. \textbf{Closing mechanism (demonstrated):} (i) $G_\lambda$ in its intrinsic Lanczos frame is a Jacobi matrix with \emph{bounded $O(1)$ coefficients}, parity-period-2; \emph{uniform boundedness verified}: $\max|A_{ij}|=4.25,4.77,5.71,4.50$ at $\lambda=9,13,17,21$ (i.e.\ $O(1)$, \emph{not} growing like $\lambda$ --- the $\mathrm{ARCH}=\mathrm{PRIME}$ balance keeps it bounded, where the raw $\sum\Lambda(n)n^{-1/2}\sim2\lambda$ diverges) --- an \emph{unconditional} property of the Weil matrix entries; (ii) Dirichlet ends (Lemma 2.3.F-Box, \emph{proved}); (iii) \emph{a bounded Jacobi with Dirichlet ends has a generic quadratic band edge}, giving $(e_k-e_0)/(e_1-e_0)\to k(k+2)$. \textbf{(iii) is demonstrated (E85, T3):} the period-2 \emph{limit} Jacobi built from the \emph{bulk-averaged} coefficients alone reproduces the shifted ladder $[0,1,2.65,4.9,7.8,11.2]\approx k(k+2)$ to ${\sim}2$--$4\%$ across $\lambda=9\dots21$ (the undershoot is the finite-chain $O((k/M)^4)$ correction $\to0$). \emph{\textbf{Why RH-neutral:}} the quadratic band edge is \emph{generic} for any bounded Jacobi --- the low modes have wavelength ${\sim}M\gg1$, so by homogenization they see only the \emph{Ces\`aro-averaged} coefficients, not the fluctuations; the ladder \emph{cannot be tuned by RH}. (This is why E84 failed --- smoothing the \emph{primes} broke the balance/marginality, the \emph{sign} --- while E85 works: smoothing the Jacobi of the already-positive $G_\lambda$, the \emph{shape}.) \textbf{Net: 2.3.F-Loc reduces to three RH-neutral, standard lemmas, none touching $\ve_0$:} (1) an unconditional bound $\sup_\lambda\max_{ij}|A_{ij}|<\infty$ (numerically verified to $\lambda=21$); (2) non-degeneracy of the band edge (quadratic, not quartic --- confirmed, the ladder is $k(k+2)$); (3) band-edge homogenization (low modes see only the Ces\`aro average --- demonstrated, standard Sturm--Liouville/effective-mass theory). A \emph{material advance} over the previous crux: 2.3.F-Loc is now \emph{standard spectral theory of a bounded Jacobi operator}, RH-neutral by construction, mechanism numerically demonstrated; the proof is the write-up of (1)--(3) --- unconditional analysis, ours to finish.
\end{remark}

\begin{theorem}[uniform gap --- PROVED modulo Lemma~\ref{lem:23F}]\label{thm:23}
By Lemma~\ref{lem:23F} and sine-diagonalization (Lemma~\ref{lem:23B}),
$\rho_\lambda^{-1}\widetilde F_{\lambda,N}\to(-\partial_\theta^2)_{\mathrm{Dir}}$, so
$\widetilde\ve_m/\widetilde\ve_1\to m^2$; subtracting the ground state and setting $m=k+1$,
$(\ve_k-\ve_0)/|\ve_0|\to(k+1)^2-1=k(k+2)$, in particular
$(\ve_1-\ve_0)/|\ve_0|\to3>0$. By Lemma~\ref{lem:23A} the limit operator is $L_\infty$
(Chebyshev). Hence $\liminf_\lambda(\ve_1-\ve_0)/|\ve_0|\ge3>0$: the uniformity of Step 2
holds. (RH-neutral: phrased as $(\ve_k-\ve_0)/|\ve_0|$, it does not decide $\operatorname{sign}\ve_0$.)
\end{theorem}
\begin{corollary}[real-rootedness in the limit, Step 3.2 --- PROVED from Thm~\ref{thm:23}]\label{cor:23G}
The uniform gap isolates the ground-state spectral projection uniformly in $\lambda$. By
Kato stability the projections converge; by Montel the $\hat\xi_\lambda$ form a normal
family with a non-degenerate limit; by Hurwitz the real-rootedness of each $\hat\xi_\lambda$
(Step 3.1, CvS) passes to the limit. Hence $\Xi$ is real-rooted on the relevant range ---
Step 3 survives $\lambda\to\infty$. \textbf{(Step 3.2 PROVED modulo Lemma~\ref{lem:23F}.)}
\end{corollary}

\subsection{Step 3 --- Reality of the zeros of $\hat\xi$ \PROVED}

\emph{Why this step, and what it receives from Step~2.} Step~2 produced, at each finite
$\lambda$, a clean ground state $\xi=\xi_\lambda$; attached to it is an entire function
$\hat\xi_\lambda$ (its Fourier transform) whose zeros are the finite-$\lambda$ analogue of the
zeta zeros. The whole reduction will compare these zeros to the true zeros of $\Xi$. For that
comparison to mean anything we must first know that, \emph{at each finite $\lambda$}, the
$\hat\xi_\lambda$ already has all its zeros real --- otherwise there would be nothing to pass to
the limit. This is the Connes--van~Suijlekom theorem, and it is essential to note that it is
\emph{not} automatic: it holds only because $A_\lambda$ is a compression of a convolution
(difference-kernel) operator, the structure Step~1 was careful to establish.

\textbf{The structural hypothesis is essential.} The naive statement ``a simple isolated even
ground state has real-rooted $\hat\xi$'' is \emph{false in general}: for any positive even
compactly supported $\xi$, the operator $I-P_\xi$ has $\xi$ as simple ground state yet $\hat\xi$
need not be real-rooted. The theorem requires the convolution / divided-difference
(Carath\'eodory--Fej\'er) structure.
\begin{theorem}[Connes--van Suijlekom 2025]
Let $Q$ be the form on $L^2[-L/2,L/2]$ associated to a real even distribution $D$ via the
Schwartz kernel $\widetilde D(x-y)$ (i.e.\ $Q$ is the compression of a convolution operator),
lower-bounded with simple isolated even spectral minimum $\xi$. Then all zeros of $\hat\xi(z)$
are real.
\end{theorem}
\begin{proof}
(i) Carath\'eodory--Fej\'er: a Hermitian PSD \emph{Toeplitz} matrix of rank $n$ has its kernel
polynomial's zeros on the unit circle --- the Toeplitz (difference-kernel) structure is exactly
what the convolution hypothesis provides. (ii) Shift $Q$ by a multiple of $\delta_0$ to make it
PSD with $\xi$ in its radical; the matrix is of the extended divided-difference class. (iii)
Finite truncations converge; Hurwitz passes reality to the limit.
\end{proof}
\noindent\textbf{Our $A_\lambda$ satisfies the hypothesis} by Lemma~1.1 (compression of
convolution by the even kernel $D_\zeta(|x-y|)$), so the theorem applies with $D=D_\zeta$.
\emph{Corollary:} $\hat\xi_\lambda$ real-rooted at each $\lambda$; RH-neutral (holds for $L_{DH}$,
also of convolution type), so non-circular.

\subsection{Step 4 --- Approximation scaffolding \PROVED}

\emph{Role.} Steps~2--3 live entirely at finite $\lambda$; to say anything about $\Xi$ (hence
about RH) we must control how well the finite-$\lambda$ world approximates the true completed
zeta function. Step~4 supplies the three quantitative estimates the limit in Step~5 consumes:
(4.1) that the truncated object $\Xi_T$ is exponentially close to $\Xi$ on the relevant window;
(4.2) that the band-limited degrees of freedom suffice to resolve every zero up to the horizon
$T^*$ (a Nyquist count, so no zero is missed); and (4.3) that the boundary/amplification losses
are negligible. None of these involve RH; they are the classical analytic bookkeeping that makes
the reduction rigorous.

\textbf{4.1.} $|\Xi_T(s)-\Xi(s)|\le C\lambda^{5/2+\delta}e^{-\pi\lambda^2}$ on $|t|\le T^*$.\quad
\textbf{4.2.} $N(T^*)=2\lambda^2\log\lambda-\lambda^2+O(\log\lambda)$ (RvM);
$\mathrm{DOF}=(\Omega/\pi)(2T^*)=4\lambda^2\log\lambda$; so
$\#\{|\gamma_\rho|\le T^*\}/\mathrm{DOF}=1-1/(2\log\lambda)<1$.\quad
\textbf{4.3.} $1-\omega\le(2/\pi)(y_0/T^*)=O(1/\lambda^2)$ (Boas); amplification
$\lambda^{|y|}\le\lambda^{1/2}$ (Phragm\'en--Lindel\"of).

\subsection{Step 5 --- The reduction (the jewel) \PROVED}

\emph{This is the step that turns RH into the sign of a number.} Everything before it was
preparation: a clean, real-rooted, well-approximating family $A_\lambda$. Step~5 now states the
payoff. Define the localized minimum
\[
\ve_\loc(\lambda)=\min_{\|f\|=1,\,f\in E_\lambda}A_\lambda(f,f) ,
\]
the smallest value the Weil form takes on bandwidth-$\lambda$ test functions. The theorem says
RH is \emph{exactly} the statement that this minimum decays slightly faster than $\lambda^{-1/2}$.
The mechanism is the one assembled in Steps~2--4: the real-rooted $\hat\xi_\lambda$ samples
$\Xi_T$ (Step~4.1), its zeros are forced onto the line by real-rootedness (Step~3), the only way
the limit can produce an off-line zero of $\Xi$ is a failure of convergence, and by Hurwitz
(Step~7) that failure is measured precisely by a non-vanishing deficit $\ve_\loc$. The uniform
gap of Step~2 (Theorem~\ref{thm:23}) is what guarantees the limit is non-degenerate, so that the
``iff'' is genuine in both directions.

\begin{theorem}
With $\ve_\loc(\lambda)=\min_{\|f\|=1,f\in E_\lambda}A_\lambda(f,f)$, one has
$\mathrm{RH}\iff\ve_\loc(\lambda)=o(\lambda^{-1/2})$.
\end{theorem}
\begin{proof}
By 4.1 approximate the zeros of $\Xi_T$; by Steps 2--3 $\hat\xi_\lambda$ is real-rooted and its
zeros sample $\Xi_T$; the deficit is $\ve_\loc$; by Hurwitz the zeros converge to those of $\Xi$
--- both real --- iff $\ve_\loc=o(\lambda^{-1/2})$.
\end{proof}

\noindent\emph{Reading.} After this theorem RH is no longer about zeros: it is the inequality
``$\ve_\loc$ is small''. Step~12.1 will sharpen ``small'' to the bare sign $\ve_0\ge0$. Steps~6--7
guard the reduction (it must reject RH-false inputs and rest on correct limiting logic), and
Steps~8--12 ask what kind of object $A_\lambda$ is and how hard that final sign really is.

\subsection{Step 6 --- Boundary fusion and the RH-false witness \PROVED}

\emph{Why a falsifiability test is needed.} A reduction of RH is worthless if it would
``prove'' RH for an object that is actually RH-false: that would mean the machinery is blind to
off-line zeros, i.e.\ circular or vacuous. Step~6 runs the only available controlled experiment.
The Davenport--Heilbronn function $\Xi_{DH}$ shares the functional equation and the same
magnitudes as $\zeta$ but is \emph{known} to have zeros off the line. We feed it through the same
construction and check that it is correctly rejected.

The band edge (carrier $\omega^*\approx\pi$) fuses with the bulk under Hurwitz. For
Davenport--Heilbronn the approximants $\hat\xi_\lambda$ are real-rooted (Step~3 applies, $L_{DH}$
being of convolution type too) but $\Xi_{DH}$ is not (off-line zero $0.8085+85.699i$); by the
contrapositive of Hurwitz the convergence $\hat\xi_\lambda\to\Xi_{DH}$ \emph{must} fail, and the
mechanism is explicit: the off-line zero adds a negative well to the symbol $a(t)$ at
$t\approx\gamma_{\rm off}$, producing $\ve_0<0$. So the form does see the violation. This both
certifies non-circularity (the test detects RH-false inputs) and localizes \emph{where} a
violation would show up (a negative well in the symbol).

\subsection{Step 7 --- Hurwitz and Laguerre--P\'olya \PROVED}

\emph{The limiting logic, made explicit.} Steps~5--6 invoked ``Hurwitz'' repeatedly; Step~7
states the two classical theorems they rest on, so the limit is on a rigorous footing. Hurwitz
controls when real-rootedness survives a limit; Laguerre--P\'olya identifies real-rootedness of
$\Xi$ with membership in the class of functions that are transforms of positive-definite forms
--- which closes the loop back to $A_\lambda\succeq0$.

Hurwitz: a uniform-on-compacts limit of real-rooted holomorphic functions is real-rooted.
$\Xi\in\mathrm{LP}\iff$ all zeros real $\iff$ RH; with Lemma~\ref{weil}, $\Xi$ is the transform
of a PSD form iff $A_\lambda\succeq0\ \forall\lambda$.

\subsection{Step 8 --- Identification of the limit operator \PROVED{} (CCM)}

\emph{From a form to an operator.} Steps~1--7 reduced RH to the sign of $\ve_0(\lambda)$ but
treated $A_\lambda$ as an abstract form. To understand the sign one wants to know \emph{what
operator} governs the bottom of the spectrum --- and here the CCM framework enters: the low
spectrum of the rescaled form $B_\lambda=\lambda^2A_\lambda$ is governed by the Slepian--Pollak
prolate operator, and the super-exponentially small eigenvalues live in its near-radical,
accessed through the conditioning map $\mathcal E$. This identification is what connects our
$A_\lambda$ to the zeta spectral triples of Steps~9--11, and it is RH-neutral: it names the
operator without deciding the sign.

$B_\lambda=\lambda^2A_\lambda$ has low spectrum governed by the Slepian--Pollak prolate operator
$W_\lambda\psi(q)=-\partial((\lambda^2-q^2)\partial)\psi+(2\pi\lambda q)^2\psi$, $|q|<\lambda$,
commuting with $P_\lambda\widehat P_\lambda P_\lambda$. The tiny eigenvalues of $A_\lambda$ (e.g.\
$2.4\cdot10^{-48}$ at $\lambda^2{=}11$) arise from the prolate near-radical via
$\mathcal E(f)(x)=x^{1/2}\sum_{n>0}f(nx)$. \textbf{8.1:} carrier $(-1)^m=\chi_m\approx(-1)^m$.
Earlier identifications \NOGO{} (\S2.6).

\subsection{Step 9 --- The semilocal prolate operator \COMP}

\emph{Building in the primes.} The bare prolate operator of Step~8 is archimedean; the zeros of
$\zeta$ require the primes. Step~9 constructs the semilocal operator $W_S$ that incorporates a
finite set $S$ of primes through its spectral measure, and --- crucially --- verifies that the
arithmetic it must carry is correct by reconstructing, as \emph{exact identities}, the
periodic-orbit (Gutzwiller) data: the prime $p$ is a closed orbit of length $\log p$ and
instability weight $p^{-1/2}$. The headline computation (9.5-bis) replaces a previously
conjectured Bessel asymptotic by the exact statement that the orbit correction
$\delta a_n$ is a matrix element of the scaling propagator $\cos(LJ)$. All of this is RH-neutral
scaffolding: it makes the operator concrete and checks it against the explicit formula, without
assuming where the zeros lie.

$W_S=scal^2+N$ on the cyclic pair $(scal,\xi_S)$; $scal=J_S$ (Jacobi), $N=\operatorname{diag}(n)$,
spectral measure
$dm_S(s)=|\Gamma_{\mathbb R}(\tfrac12+is)|^2\prod_{p\in S}|1-p^{-(1/2+is)}|^{-2}ds$. Even
$\Rightarrow b_n=0$, $W_S=J_S^2+\operatorname{diag}(n)$. \emph{The object CCM deferred.}

\textbf{9.1.} $S=\{\infty\}$: $a_n=\sqrt{(n-\tfrac12)n}$ reproduces CCM exactly to $n=15$.
\textbf{9.2.} $|\delta a_n/a_n|\approx C'p^{-1/2}(n\log p)^{-1/2}$, $C'\approx0.85$, across
$p=2,\dots,23$: $p^{-1/2}$ is the Gutzwiller orbit weight, $(n\log p)^{-1/2}$ the Bessel envelope
$J_2(2n\log p)$. \emph{Linear-response identity (validated):} with $w_\infty=|\Gamma_{\mathbb
R}(\tfrac12+is)|^2$ and the archimedean (Meixner--Pollaczek) orthonormal $p_n$,
$\delta\log a_n=\sum_{p,k}\tfrac{p^{-k/2}}{k}I_n(k\log p)$,
$I_n(L)=\int_{\mathbb R}\cos(Ls)(p_n^2-p_{n-1}^2)w_\infty\,ds$; verified vs.\ direct $\delta a_n$ for
$p=23,29,31,37$ (ratio $0.997,0.997,1.009,0.995$, std $\le0.05$).
\textbf{9.5-bis (exact closed form, PROVED).} Since $w_\infty$ is the spectral measure of the
Jacobi operator $J$ ($a_n=\sqrt{(n-\tfrac12)n}$, $b_n=0$), one has
$\int e^{iLs}p_np_m w_\infty\,ds=\langle p_n,e^{iLJ}p_m\rangle=(e^{iLJ})_{nm}$, hence
$\boxed{I_n(L)=[\cos(LJ)]_{n,n}-[\cos(LJ)]_{n-1,n-1}}$. This is exact (verified to $10^{-16}$ for
$L=\log2,\log3$, all $n\le14$). Thus
$\delta\log a_n=\sum_{p,k}\tfrac{p^{-k/2}}{k}\bigl([\cos(k\log p\,J)]_{n,n}-[\cos(k\log
p\,J)]_{n-1,n-1}\bigr)$: $\delta a_n$ is the diagonal of the scaling-operator propagator
$\cos(LJ)$, $L=\log p^k$ the orbit length --- the Gutzwiller picture as an exact identity.
\textbf{(PROVED, RH-neutral)}

\medskip\noindent\textbf{9.5-ter --- Fine oscillation law of $I_n(L)$ \PROVED.}
\emph{The exact $n\to\infty$ asymptotics, closing the fine-oscillation law and superseding the
conjectured Bessel form $C_0J_2(2nL)$.} The matrix $J$ ($a_n=\sqrt{n(n-\tfrac12)}$, $b_n=0$)
is the Jacobi matrix of the orthonormal symmetric Meixner--Pollaczek polynomials with
$\lambda=\tfrac14$ and spectral measure $w_\infty(s)=|\Gamma_{\mathbb R}(\tfrac12+is)|^2$;
equivalently $J=2K_1$ is twice the boost generator of the $su(1,1)$ discrete series $D_{1/4}$
(from $a_n=\sqrt{n(n+2k-1)}$ with $k=\tfrac14$).
\begin{theorem}[fine oscillation law]\label{thm:95}
Fix $L>0$ and set $\omega(L):=2\arcsin(\tanh L)=2\,\mathrm{gd}(L)$ (twice the Gudermannian;
$\sin(\omega/2)=\tanh L$, $\cos(\omega/2)=\mathrm{sech}\,L$). Then as $n\to\infty$,
\[
[\cos(LJ)]_{n,n}=(\pi n\sinh L)^{-1/2}\cos\!\big((n+\tfrac14)\omega(L)-\tfrac\pi4\big)+O(n^{-3/2}),
\]
\[
I_n(L)=\frac{2\tanh L}{\sqrt{\pi\sinh L}}\;n^{-1/2}\cos\!\big((n-\tfrac14)\omega(L)+\tfrac\pi4\big)+O(n^{-3/2}).
\]
\end{theorem}
\begin{proof}
(i) By the spectral theorem $[\cos(LJ)]_{n,n}=\int_{\mathbb R}\cos(Ls)p_n(s)^2w_\infty(s)\,ds$,
real since $b_n=0$ makes $p_n^2w_\infty$ even. (ii) In the bulk $|s|<2a_n$ the Plancherel--Rotach
law (Li--Wong; Liouville--Green for the recurrence) gives
$p_n\sqrt{w_\infty}=\sqrt{2/(\pi R_n)}\cos\Psi_n+O(n^{-1})$, $R_n=\sqrt{4a_n^2-s^2}$,
$\Psi_n(s)=\sum_{k\le n}\arccos(s/2a_k)$, so $p_n^2w_\infty=(\pi R_n)^{-1}(1+\cos2\Psi_n)+\dots$
(iii) In $\int\cos(Ls)\cos2\Psi_n\,(\pi R_n)^{-1}ds$ the stationary points solve $2\Psi_n'(s)=\mp L$;
with $\Psi_n'(s)=-\sum_k(4a_k^2-s^2)^{-1/2}\approx-\tfrac12\mathrm{arccosh}(2n/|s|)$ this gives
$|s_*|=2n\,\mathrm{sech}\,L$, and by the envelope theorem the $n$-frequency is
$\omega=\tfrac{d}{dn}2\Psi_n(s_*)=2\arccos(s_*/2a_n)=2\arccos(\mathrm{sech}\,L)=2\arcsin(\tanh L)$
(using $1-\mathrm{sech}^2L=\tanh^2L$) --- exactly the dispersion, $\neq2L$, which is why the Bessel
form failed. (iv) At $s_*$, $R_n(s_*)=2n\tanh L$ and the shift $a_k=k-\tfrac14+O(1/k)$ gives degree
$\nu=n-\tfrac14$; equivalently $[\cos(LJ)]_{n,n}=\mathrm{sech}\,L\,P_{n-1/4}(\cos\omega)+O(n^{-3/2})$,
and Laplace--Heine $P_\nu(\cos\theta)=\sqrt{2/(\pi\nu\sin\theta)}\cos((\nu+\tfrac12)\theta-\tfrac\pi4)+O(\nu^{-3/2})$
with $\sin\omega=2\tanh L\,\mathrm{sech}\,L$ yields the stated amplitude $(\pi n\sinh L)^{-1/2}$ and
phase. (v) $\cos\Phi_n-\cos\Phi_{n-1}=-2\sin(\omega/2)\sin(\Phi_n-\omega/2)$ with
$\sin(\omega/2)=\tanh L$ gives $I_n$.
\end{proof}
\noindent\emph{Validation (exact engine, $N\gg n$):} the law matches the computed
$[\cos(LJ)]_{n,n}$ and $I_n$ to relative error $\le3\cdot10^{-4}$ (decreasing as $O(n^{-1})$),
uniformly over orbit lengths $L=\log2,\log3,\log4,\log5,\log8$, $n\le2500$; the fitted dispersion
matches $2\arcsin(\tanh L)$ to 4 digits on $L\in[0.3,2.2]$ and $B_n\sqrt{\pi n\sinh L}=1$ to 4
digits independently of $L$. This \textbf{refutes $C_0J_2(2nL)$ for a precise reason} (the true
frequency is the Gudermannian $2\arcsin(\tanh L)$, not $2L$; the earlier ``envelope grows''
refutation observed the pre-asymptotic $n\le80$ regime). The Gudermannian dispersion is the
$su(1,1)$ signature: the circular angle $\omega$ is conjugate to the hyperbolic orbit rapidity $L$.
\textbf{(PROVED --- Plancherel--Rotach $+$ stationary phase / Laplace--Heine; constants verified
to 4 digits.)}
\textbf{9.3.} Primes compress the moments: $\mu_2/\mu_0=0.50\to0.16\to0.08\to0.05$.
\textbf{9.4.} $a_n^S\sim n\Rightarrow$ Carleman diverges $\Rightarrow J_S$ essentially
self-adjoint; the negative eigenvalues (zeros) live in the orthogonal complement (Sonin).
\PROVED

\subsection{Step 10 --- Counting and periodic orbits \PROVED/\COMP}

\emph{Confirming the operator is the right one.} An operator proposed to carry the zeros must
reproduce the zero-counting function of $\zeta$, smooth part and fluctuation alike. Step~10 is
that consistency check: the smooth count matches the CCM semiclassical law, and --- the decisive
point --- the fluctuation $S(E)=\frac1\pi\arg\zeta(\tfrac12+iE)$ is reproduced \emph{exactly} as
a Gutzwiller sum over prime orbits, numerically to three digits. Together with Step~9 this fixes
the arithmetic of the operator beyond reasonable doubt; what remains undecided is only the one
thing that is RH (the sign / the reality of the negative spectrum), handled in Steps~11--12.

\textbf{10.1.} $N(E)=N_{\rm smooth}(E)+S(E)$,
$N_{\rm smooth}=\tfrac{E}{2\pi}(\log\tfrac{E}{2\pi}-1)+\tfrac78$,
$S(E)=\tfrac1\pi\arg\zeta(\tfrac12+iE)$; CCM's $\sigma(E,\lambda)$ reproduces $N_{\rm smooth}$.
\textbf{10.2.} $S(E)=-\tfrac1\pi\sum_{p,k}(kp^{k/2})^{-1}\sin(kE\log p)$; verified
$S_{\rm real}=S_{\rm primes}$ ($-0.293$ vs $-0.294$ at $E{=}40$). A Gutzwiller sum: orbit $=p$,
length $\log p$, weight $p^{-1/2}$, repetitions $k$. \PROVED
\textbf{10.3.} $H_\lambda=(p^2-\lambda^2)(q^2-\lambda^2)$ has open (scattering) orbits $\Rightarrow$
only $N_{\rm smooth}$; the closed orbits $\log p$ come from the finite places.

\subsection{Step 11 --- Sonin space and de Branges function \OPEN$\to$\WALL}

\emph{Where the zeros actually live, and why this is the wall.} The positive spectrum built in
Steps~8--10 is archimedean scaffolding; the nontrivial zeros sit in the \emph{negative} spectrum
of the self-adjoint prolate operator, carried by the Sonin space. Step~11 constructs this
negative spectrum (rigorously, via a monotone Pr\"ufer flow) and validates its counting against
the semiclassical law, then asks the RH question in de Branges language: are the associated
$\gamma_n$ real? Every attempt to force this by a \emph{positivity} (Hermite--Biehler, a scalar
de Branges quotient, a positive Hamiltonian $H_S\succeq0$) is refuted with its reason --- the
scalar positivity is false for $\zeta$ (Conrey--Li/Sarnak), and the correct object is an
indefinite, finite-Pontryagin-index transfer matrix. The reality of the $\gamma_n$ is reached
here from yet another independent direction and remains undecided: \emph{this is the same wall},
now wearing its de Branges costume.

$S_\lambda=\{f\in L^2(\mathbb R)^{ev}:f=0,\hat f=0\text{ on }[-\lambda,\lambda]\}$ carries the
negative spectrum of $W_{sa}$ (zeros, $\sim-\gamma_n^2$ UV).
\textbf{11.1.} $E_S(z)=1/\Lambda_S(\tfrac12-iz)$ entire, $|E_S(x)|^{-2}=dm_S(x)$ (verified).
\textbf{11.2.} $E_S$ \emph{not} Hermite--Biehler: at $z=2+i$, $|E_S(z)|=5.13<|E_S(\bar z)|=5.80$
(gamma factor decays in the upper half-plane). \NOGO
\textbf{11.2-bis (structural no-go).} The scalar
$\Theta_S(z)=\Lambda_S(\tfrac12-iz)/\Lambda_S(\tfrac12+iz)$ has poles/zeros at
$z=2\pi m/\log p\pm i/2\in\mathbb C_+$, hence is not Schur/contractive in $\mathbb C_+$: it cannot
be a characteristic function of a self-adjoint extension nor a de Branges quotient. Any scalar
Robin condition from $\Lambda_S$ inherits these singularities, and (DH sharing the partial-Euler
structure) is exactly what marginality (12.2) rules out. \emph{The right finite-$S$ object is a
$2\times2$ $J$-unitary transfer matrix} (Tate), giving a de Branges canonical system $J\dot
Y=zH_SY$; the wall is $H_S\succeq0$. \NOGO\ (sharpened $\to$ \S4.B)
\textbf{11.3.} HB $=$ reality of zeros requires $s\leftrightarrow1-s$. $\Lambda(s)=\Lambda(1-s)$
makes $\xi(\tfrac12+iz)$ real, even, zeros $\gamma_n$ (verified
$\operatorname{Im}\xi(\tfrac12+it)\sim10^{-23}$). The partial $\Lambda_S$ has no functional
equation: its zeros are trivial (neg.\ imag.\ axis) and Euler ($\operatorname{Im}z=-\tfrac12$),
not $\gamma_n$. \PROVED
\textbf{11.3-bis (Sonin negative spectrum --- constructed and validated).} By CCM
(Thm 2.6.iv) the negative spectrum of $W_{\mathrm{sa}}$ is that of
$W_\lambda\psi=-\partial((\lambda^2-q^2)\partial)\psi+(2\pi\lambda q)^2\psi$ on $L^2(\lambda,\infty)$
with BC (16) regular at $q=\lambda$ and (17) at $q=\infty$. With $u=R\sin\theta$,
$(q^2-\lambda^2)u'=R\cos\theta$, the equation
$((q^2-\lambda^2)u')'+((2\pi\lambda q)^2-\xi)u=0$ gives the \emph{monotone} Pr\"ufer flow
$\theta'=\cos^2\theta/(q^2-\lambda^2)+((2\pi\lambda q)^2-\xi)\sin^2\theta$, $\theta(\lambda)=\pi/2$.
For $\xi<0$ both terms are positive, so $\theta$ increases in $q$ and in $-\xi$: every eigenvalue
is a crossing $\theta=m\pi$, none skipped or spurious --- reliable even at the limit-circle
oscillatory endpoint where the asymptotic functional (17) (valid only for $x\gg\mu/12$) fails.
\emph{Validation:} the counting $N(E)=\#\{-\xi\le E^2\}$ has spectral density $dN/dE$ matching the
CCM law $d(2\sigma)/dE$ to $\approx0.9$ uniformly across $\lambda=1,2,3$ (increments
$\Delta N/\Delta(2\sigma)=0.85$--$0.96$, no $\lambda$-bias); the cumulative count correctly gives
zero eigenvalues below the first ($\lambda=3$: first at $E\approx16$). \PROVED\ (reliable
construction; density validated)
\textbf{11.4.} Closing HB with the functional equation $=$ $\gamma_n$ real. \WALL

\subsection{Step 12 --- The wall \WALL}

\emph{The single statement that is RH.} Step~12 names the wall in its cleanest form and proves
the one structural fact that explains why sixty phases failed to climb it. By 12.1 the entire
chain collapses to the positivity $(P):\ve_0(\lambda)\ge0$ for all $\lambda$, which is RH. The
marginality theorem (12.2) then shows why $(P)$ is so hard: the negative central well of the
symbol and the band resolution are in a fixed, scale-invariant ratio $\sqrt2/\pi$ that depends
only on the \emph{magnitudes} $|\Lambda(n)|$ --- and Davenport--Heilbronn has the very same
magnitudes yet $\ve_0<0$. Hence no magnitude/size method can ever decide the sign of $\ve_0$:
the sign is a \emph{phase} invariant. This is the quantitative core of the program's negative
discovery, and it tells any future prover exactly what kind of argument is required --- one that
is phase-sensitive and encodes the location of the zeros by construction, which is precisely the
CCM operator route whose stability rests on the new Lemma~\ref{lem:23F}.

\textbf{12.1.} $(P):\ve_0(\lambda)\ge0\ \forall\lambda\iff$ RH (Lemma~\ref{weil} and DH).
\textbf{12.2 (marginality).} Central well half-width $\delta_\lambda=(2|a(0)|/a''(0))^{1/2}$,
resolution $r_\lambda=\pi/\Omega$; with $P_\lambda\sim2\lambda$,
$M_2\sim8\lambda\log^2\lambda$ (PNT), $a(0)\sim-2\lambda$, $a''(0)\sim M_2$:
$\frac{2\delta_\lambda}{r_\lambda}=\frac{\sqrt2}{\pi}\approx0.45$. Since $\delta_\lambda,r_\lambda$
depend only on $\{|\Lambda(n)|\}$, no magnitude functional decides $\operatorname{sign}\ve_0$ (DH,
same magnitudes, $\ve_0<0$). The sign is a phase invariant --- the quantitative form of MW-4
(wrong sign) and MW-2 (phase unreachable by magnitude). \PROVED

\section{Future work}
\paragraph{4.A --- Closed.} The law of $\delta a_n$ (Step 9.5-bis) is now exact:
$I_n(L)=[\cos(LJ)]_{n,n}-[\cos(LJ)]_{n-1,n-1}$, a matrix element of the scaling-operator
propagator, superseding the (numerically refuted) Bessel/DIK conjecture. No RH-neutral residue
remains in Step 9.
\paragraph{4.B --- The frontier: adelic transfer matrices and canonical-system positivity.}
By 11.2-bis the right finite-$S$ object is not the scalar partial product but a $2\times2$
$J$-unitary transfer matrix from Tate's local functional equation: take $\varphi=\varphi_\infty
\otimes\bigotimes_p\varphi_p$ on $\mathbb A_{\mathbb Q}$, $\varphi_p=\mathbf 1_{\mathbb Z_p}$ for
$p\in S$, and keep both $Z_p(\varphi_p,s)$ and the dual $Z_p(\widehat\varphi_p,1-s)$ as entries of
$T_p(z)\in SU(1,1)$. Then $T_S=T_\infty\prod_{p\in S}T_p$ is the transfer matrix of a de Branges
canonical system $J\dot Y=zH_SY$ (this is why scalar HB fails: for $p\notin S$,
$\widehat\varphi_p\ne\varphi_p$, and the dual side restores the discarded data). Three linked
targets:
(1) \emph{Boundary triple (RH-neutral, doable):} $W_\lambda$ on $|q|<\lambda$ has limit-circle log
endpoints, deficiency $(4,4)$, $2$ even; tabulate the matrix Weyl function $M_\lambda(z)$,
$\operatorname{Spec}_{\rm disc}(W_{\lambda,B})=\{E:\det(B-M_\lambda(E))=0\}$; classify the
Fourier-invariant $B$ and compare $B_{\rm CCM}$ with Burnol's $K_\lambda/Z_\lambda$.
(2) \emph{The Hamiltonian:} prove $T_p$ $J$-unitary on $\mathbb R$, $J$-contractive in $\mathbb
C_+$, so $H_S\succeq0$ and $E_S$ is HB; Hurwitz closes $S\to\mathbb P$. \emph{The wall is
$H_S\succeq0$} (compare carefully with Conrey--Li IMRN 2000).
(3) \emph{Dynamical symmetry:} $s\leftrightarrow1-s$ emerges from the $\widehat{\mathbb Z}^\times$
symmetry of Bost--Connes KMS states at $\beta=1$; the semilocal Sonin operator should be a
restriction of a Bost--Connes-type Hamiltonian. \emph{Decisive RH-neutral test:} can the
boundary condition at $q=\pm\lambda$ be stated without reference to the zeros? If not, that is
where RH is inserted.
\paragraph{4.C --- The wall.} $(P):\ve_0\ge0$ is RH (MW-1), immune to magnitude methods (Step
12.2 $=$ MW-4/MW-2 quantitative). The collapse theorem says any unconditional route must be
phase-sensitive and encode the zeros by construction --- the CCM/Hilbert--P\'olya route of 4.B,
the only non-circular candidate, not closed.

\paragraph{Summary of status.} Steps 1--10 are proved or constructed-and-validated. The
sixty-phase unconditional results: the localized Weil detector (Arc B, Fronts A, B); the
reduction $\mathrm{RH}\iff\ve_\loc=o(\lambda^{-1/2})$; the structure theorem; real-rootedness
(CvS); the approximation scaffolding; the marginality theorem; and the construction and
validation of the semilocal prolate operator with its periodic-orbit weight. Step 9.5 is now
proved (Theorem~\ref{thm:95}, the Gudermannian fine oscillation law); Step 11 is open and,
with Step 12, is the wall, reached identically from primes, zeros,
operator, de Branges, and the seven structural walls of the earlier arcs: the reality of the
nontrivial zeros, i.e.\ RH. The contribution is a rigorous localization of RH onto a single
phase-sensitive statement, the elimination --- with precise obstructions --- of every
alternative the program could construct or find in the literature, the CCM scaffolding around
the wall, and the precise identification of the missing ingredient.

\begin{thebibliography}{9}
\bibitem{CCM} A.~Connes, C.~Consani, H.~Moscovici, \emph{Zeta spectral triples}, arXiv:2511.22755 (2025).
\bibitem{CvS} A.~Connes, W.~D.~van Suijlekom, \emph{Quadratic Forms, Real Zeros and Echoes of the Spectral Action}, arXiv:2511.23257 (2025).
\bibitem{CCzc} A.~Connes, C.~Consani, \emph{Spectral triples and $\zeta$-cycles}, arXiv:2106.01715 (2021).
\bibitem{CMpro} A.~Connes, H.~Moscovici, \emph{Prolate spheroidal operator and Zeta}, arXiv:2112.05500 (2021).
\bibitem{CMsemi} A.~Connes, H.~Moscovici, \emph{Zeta zeros and prolate wave operators}, arXiv:2310.18423 (2024).
\bibitem{CCweil} A.~Connes, C.~Consani, \emph{Weil positivity and trace formula, the archimedean place}, Selecta Math. 27 (2021).
\bibitem{KS} P.~Kurasov, P.~Sarnak, \emph{Stable polynomials and crystalline measures} (2020).
\bibitem{Mont} H.~L.~Montgomery, \emph{The pair correlation of zeros of the zeta function} (1973).
\end{thebibliography}

\end{document}
